\chapter{Synthetic test bench}
\label{ch:synthetic}

\begin{revised}[title={New chapter, added 14 October 2026}]
This chapter documents \file{scripts/sysid/simulation/synth_bag.py} and its tests, which were
added after the first edition of these notes. All numbers below are produced by
\file{docs/sysid_notes/gen_synthdata.py}, which runs the repository's own propagation and fit
code on freshly generated bags. \Cref{app:changelog} lists every change made in this revision.
\end{revised}

\section{Why ground truth is needed}

A fit of a real recording returns numbers, but nothing says whether they are \emph{right}: the
true coefficients of POLARIS are exactly what is unknown. An estimator can be biased by its
own construction (filtering, delays, the choice of regressors), and that bias is invisible in
real data. It looks like a property of the vehicle.

The standard remedy is to test the method on data whose truth is known
\cite{ljung1999,klein2006}. A simulated vehicle with chosen coefficients is driven through the
same manoeuvres, its sensors are imitated with chosen delays and noise, and the result is
written in the same format as a real recording. Any difference between the fitted and the
chosen coefficients is then an error \emph{of the method}, and it can be traced to the
imperfection that caused it by switching imperfections on one at a time.

\begin{keyidea}
Passing on synthetic data is \emph{necessary}, not sufficient. A fit that cannot recover the
coefficients of a bag whose truth is known will not recover them from the lake either. A fit
that does recover them may still fail on effects the simulation leaves out (\cref{sec:synth-limits}).
\end{keyidea}

\section{What is simulated}

\file{synth_bag.py} integrates the surge model of \cref{ch:surge} on a \qty{1}{\kilo\hertz} grid
with the coefficients and thrust map of a \emph{truth model} file (by default the dummy
\file{dummy_bluerov2_heavy_t500.yaml}: BlueROV2 coefficients with the real T500 map):
\begin{equation}
  \mtot\,\dot u = k\,T_\mathrm{act}(t) + F_\mathrm{off} - X_u u - \Xuu|u|u .
  \label{eq:synth-plant}
\end{equation}
The imperfections are layered on top as follows.
\begin{description}[style=nextline]
  \item[Actuation] The force follows the \emph{logged} servo PWM $p_\mathrm{log}$ after a pure delay $L$
        (\code{actuation_lag_s}) and a first-order motor spin-up $\tau_m$
        (\code{spinup_tau_s}):
        \begin{equation}
          \tau_m\,\dot T_\mathrm{act} = F\bigl(p_\mathrm{log}(t - L),\, V(t)\bigr) - T_\mathrm{act}.
        \end{equation}
        $k$ (\code{thrust_scale}) scales the true thrust against the table, which models a
        datasheet error. $F_\mathrm{off}$ (\code{offset_n}) is a constant tether or current force.
  \item[Battery] $V(t)$ sags linearly from \code{v_start} to \code{v_end}. The BMS publishes it at
        \qty{1}{\hertz} as the same text status (\code{"16.20V"}) as \code{bms_node}.
  \item[EKF] The surge velocity passes through a first-order lag $\tau_e$
        (\code{ekf_lag_s}), $\tau_e\dot{\hat u} = u - \hat u$, plus white noise. The pose is the
        integral of $\hat u$ along a constant heading, so pose and twist stay consistent.
  \item[IMU] Specific force in FLU, $\vect{f} = \dot{\vect{v}} - \vect{R}^\top[0,0,-g]^\top$, which at
        constant trim pitch $\theta_0$ gives $f_x = \dot u - g\sin\theta_0$ and
        $f_z = g\cos\theta_0$ (\cref{ch:timing}), plus white noise.
  \item[DVL] The body velocity rotated into \code{dvl_a50_link} with the real mount from
        \code{/tf_static} ($\vect{v}_\mathrm{dvl} = \vect{R}_\mathrm{d}^\top\vect{v}_b$), optionally
        delayed, plus white noise.
  \item[Scripts] The \code{/pixhawk/manual_control} messages and the \code{/sysid/event}
        markers are reproduced message for message from \file{_step_common.py}: 20 Hz commands,
        \code{thrust_on}, \code{thrust_cut}, a neutral burst and \code{done}.
\end{description}

\begin{table}[htbp]
\centering
\caption{Imperfections of the synthetic vehicle: command-line flag, default value, and
value with \code{--ideal}. Defaults are plausible guesses, not measurements.}
\label{tab:synth-config}
\small
\begin{tabular}{@{}l l S[table-format=2.3] S[table-format=2.3]@{}}
\toprule
Quantity & Flag & {Default} & {Ideal} \\
\midrule
Actuation lag $L$ / \unit{\second} & \code{--actuation-lag-s} & 0.15 & 0 \\
Spin-up $\tau_m$ / \unit{\second} & \code{--spinup-tau-s} & 0.05 & 0 \\
Thrust scale $k$ & \code{--thrust-scale} & 1 & 1 \\
Force offset / \unit{\newton} & \code{--offset-n} & 0 & 0 \\
Battery start, end / \unit{\volt} & \code{--v-start}, \code{--v-end} & {16.2, 15.6} & {16.2, 16.2} \\
EKF lag $\tau_e$ / \unit{\second} & \code{--ekf-lag-s} & 0.10 & 0 \\
EKF noise / \unit{\meter\per\second} & \code{--ekf-noise-mps} & 0.01 & 0 \\
IMU noise / \unit{\meter\per\second\squared} & \code{--imu-noise-mps2} & 0.02 & 0 \\
DVL noise / \unit{\meter\per\second} & \code{--dvl-noise-mps} & 0.005 & 0 \\
DVL delay / \unit{\second} & \code{--dvl-lag-s} & 0 & 0 \\
PWM per command unit / \unit{\micro\second} & \code{--pwm-per-cmd} & 0.4 & 0.4 \\
Trim pitch $\theta_0$ / \unit{\radian} & \code{--trim-pitch-rad} & 0.02 & 0.02 \\
\bottomrule
\end{tabular}
\end{table}
\FloatBarrier

The PWM mapping $p = 1500 + 0.4\,c$ (command 500 $\to$ \qty{1700}{\micro\second}) is an
\emph{assumption}. ArduSub's real mapping depends on its pilot gain and motor parameters. The
first lake bag that records both \code{/pixhawk/manual_control} and
\code{/pixhawk/servo_output_raw} settles it.

\begin{table}[htbp]
\centering
\caption{Topics written by \file{synth_bag.py}, with message type and rate. Types and frames are
identical to the real recordings (ROS~2 Jazzy).}
\label{tab:synth-topics}
\small
\begin{tabularx}{\textwidth}{@{}l l l X@{}}
\toprule
Topic & Type & Rate & Note \\
\midrule
\code{/odometry/filtered/local} & \code{nav_msgs/Odometry} & \qty{30}{\hertz} & \code{odom} $\to$ \code{base_link} \\
\code{/pixhawk/servo_output_raw} & \code{std_msgs/Int16MultiArray} & \qty{10}{\hertz} & no header, other channels \num{1500} \\
\code{/pixhawk/manual_control} & \code{std_msgs/Int16MultiArray} & \qty{20}{\hertz} & only while a script runs \\
\code{/imu/data} & \code{sensor_msgs/Imu} & \qty{100}{\hertz} & \code{imu_link} \\
\code{/sensors/dvl/velocity} & \code{marine_acoustic_msgs/Dvl} & \qty{7.6}{\hertz} & definition in \file{simulation/Dvl.msg} \\
\code{/diagnostics} & \code{diagnostic_msgs/DiagnosticArray} & \qty{1}{\hertz} & BMS current, voltage, percentage \\
\code{/sysid/event} & \code{std_msgs/String} & events & JSON as \file{_step_common.py} \\
\code{/sysid/synthetic_truth} & \code{std_msgs/String} & once & the truth as JSON \\
\code{/tf_static} & \code{tf2_msgs/TFMessage} & once & \code{imu_link}, \code{dvl_a50_link} \\
\bottomrule
\end{tabularx}
\end{table}
\FloatBarrier

Next to every bag, \file{<stem>_truth.yaml} records the true $m$, $\Xud$, $\mtot$, $X_u$,
$\Xuu$, every timing parameter (including the equivalent \code{servo_delay_s}, see below),
the full configuration, the manoeuvre programme and the time of every event.

\section{Using it}

\begin{lstlisting}[language=bash,caption={Generating synthetic bags (run from the repository root).}]
# the full suite (four steps and a coast-down) with default imperfections
python scripts/sysid/simulation/synth_bag.py --program suite

# perfect data: no delay, no noise, constant voltage
python scripts/sysid/simulation/synth_bag.py --program steps --ideal

# perfect data except one imperfection, here a 0.2 s actuation lag
python scripts/sysid/simulation/synth_bag.py --program steps --ideal --actuation-lag-s 0.2

# several noise realisations of the same experiment
for s in 0 1 2 3 4; do python scripts/sysid/simulation/synth_bag.py --seed $s; done
\end{lstlisting}

Output goes to \file{recordings/sysid/synthetic/} (git-ignored) unless \code{--out} is given.
A suggested workflow for each fit script:
\begin{enumerate}
  \item \textbf{Ideal bag.} The fit must return the truth up to discretisation. If it does not,
        the method itself is biased, and the remaining steps would only add noise on top.
  \item \textbf{One imperfection at a time.} Each bias can then be attributed to its cause
        (\cref{tab:synth-fit}).
  \item \textbf{Default imperfections, several seeds.} The spread of the estimates over seeds is the
        statistical uncertainty of the method. Their mean against the truth is its bias.
  \item \textbf{Only then real bags.}
\end{enumerate}

\section{Result 1: the propagation reproduces the truth}

\Cref{tab:synth-prop} runs \code{propagate()} from \file{propagate_model.py} with the truth
model over two bags of the \code{steps} programme: one ideal, one with only a
\qty{0.2}{\second} actuation lag.

\begin{table}[htbp]
\centering
\caption{RMS surge error of the free-run propagation of the truth model over synthetic bags,
for two actuation lags $L$ and three values of \code{thrust.servo_delay_s}.}
\label{tab:synth-prop}
\begin{tabular}{@{}c c l@{}}
\toprule
Lag $L$ / \unit{\second} & \code{servo_delay_s} / \unit{\second} & Surge RMS / \unit{\meter\per\second} \\
\midrule
\inputrows{data/synth_prop_rows.tex}
\bottomrule
\end{tabular}
\end{table}
\FloatBarrier

On the ideal bag the propagation matches the truth to about \qty{1e-8}{\meter\per\second}. This
single number checks every link at once: bag format and reading, timestamps, BMS voltage
parsing, the T500 map, zero-order hold and RK4.

\begin{caution}[title={The sign of \code{servo_delay_s}}]
The lagged bag is reproduced exactly with $\code{servo_delay_s} = -0.2$, and a \emph{positive}
value makes it \emph{worse}. Both \file{propagate_model.py} and \file{coupling_check.py}
\emph{subtract} \code{servo_delay_s} from the servo log time, so a positive value moves the
thrust \emph{earlier} than its log entry. That compensates transport delay (the log arrives after
the Pixhawk applied the PWM) but not motor spin-up (the thrust arrives after the log). If the IMU
shows the vehicle accelerating \qty{150}{\milli\second} after the logged PWM step,
\textbf{the value to enter is $\code{servo_delay_s} = -0.15$}.
\file{maneuvers.md} and the first edition of \cref{ch:timing} read as if it were positive. The
convention is unchanged in the code. Whether to keep it and document it, or to flip it, is a
decision still to be made (\cref{ch:status}).
\end{caution}

\begin{figure}[htbp]
\centering
\begin{tikzpicture}
\begin{axis}[wide,xlabel={Time after \code{thrust_on} / \unit{\second}},ylabel={$u$ / \unit{\meter\per\second}},
  xmin=-0.5,xmax=2.5,ymin=-0.03]
\addplot[black,very thick] table[x=t,y=truth]{data/synth_step.dat};
\addplot[gray,only marks,mark=*,mark size=0.8pt] table[x=t,y=ekf]{data/synth_step.dat};
\addplot[red,dashed] table[x=t,y=delay0]{data/synth_step.dat};
\addplot[blue,dashdotted] table[x=t,y=delaycomp]{data/synth_step.dat};
\legend{truth,EKF,model $\code{servo_delay_s}=0$,model $\code{servo_delay_s}=-0.15$}
\end{axis}
\end{tikzpicture}
\caption{Surge speed around the surge-500 step of a synthetic suite bag with default
imperfections: true speed, EKF measurement, and free-run propagation of the truth model with
\code{servo_delay_s} of \num{0} and \qty{-0.15}{\second}.}
\label{fig:synth-step}
\end{figure}
\FloatBarrier

\Cref{fig:synth-step} shows the same on a bag with all default imperfections. Without
compensation the model leads the truth by the actuation lag. With $-0.15$ it follows the
truth closely. The small remaining lead is the \qty{0.05}{\second} motor spin-up, which a pure
time shift cannot represent. The EKF trails both, because of its own
\qty{0.1}{\second} lag, which is the filter lag of \cref{ch:timing}.

\section{Result 2: biases of the equation-error fit}

\Cref{tab:synth-fit} applies the own-axis least-squares fit of \file{coupling_check.py}
(\cref{ch:estimation}; Savitzky--Golay window \qty{1}{\second}, \qty{50}{\hertz} grid) to suite
bags that each carry a single known imperfection. The truth is $\mtot = \qty{19.86}{\kilogram}$,
$X_u = \qty{13.7}{\newton\second\per\meter}$, $\Xuu = \qty{141}{\newton\second\squared\per\meter\squared}$.

\begin{table}[htbp]
\centering
\caption{Coefficients returned by the equation-error fit of \file{coupling_check.py} on
synthetic suite bags with one imperfection each, with the deviation from the truth in brackets.}
\label{tab:synth-fit}
\small
\begin{tabular}{@{}l l l l@{}}
\toprule
Bag & $\mtot$ / \unit{\kilogram} & $X_u$ / \unit{\newton\second\per\meter} & $\Xuu$ / \unit{\newton\second\squared\per\meter\squared} \\
\midrule
\inputrows{data/synth_fit_rows.tex}
\bottomrule
\end{tabular}
\end{table}
\FloatBarrier

Four lessons follow, each one invisible on real data.
\begin{enumerate}
  \item \textbf{The ideal bag is not fitted exactly.} $\mtot$ comes out \qty{14}{\percent} high and
        $X_u$ \qty{18}{\percent} high on perfect data. The cause is the \qty{1}{\second}
        smoothing window, which is long against $\taueff \approx \qty{0.2}{\second}$
        (\cref{sec:taueff}). Smoothing $u$ and then forming $|u|u$ is not the same as smoothing
        $|u|u$, and the step edges are smeared over the whole window.
  \item \textbf{Delays inflate the inertia}, exactly as argued in \cref{ch:timing}. A
        \qty{0.15}{\second} actuation lag raises the error of $\mtot$ from \qty{14}{\percent} to
        \qty{65}{\percent}, and a \qty{0.1}{\second} EKF lag raises it to \qty{53}{\percent}. Drag barely moves.
  \item \textbf{A thrust-scale error rescales everything} by $1/k$, as derived in
        \cref{sec:identifiability}. Here $\mtot$ happens to land on the truth, because the
        \qty{14}{\percent} smoothing bias and the factor $1/1.15$ cancel. A single coefficient
        that ``looks right'' proves nothing.
  \item \textbf{Errors add up.} With all defaults combined, $\mtot$ is twice its true value.
\end{enumerate}

\paragraph{The smoothing window.} \Cref{tab:synth-window} repeats the fit for shorter windows.

\begin{table}[htbp]
\centering
\caption{Total inertia $\mtot$ returned by the equation-error fit for four Savitzky--Golay window
lengths, on the ideal, the lagged and the default synthetic suite bags (truth \qty{19.86}{\kilogram}).}
\label{tab:synth-window}
\small
\begin{tabular}{@{}c l l l@{}}
\toprule
Window / \unit{\second} & ideal & lag \qty{0.15}{\second} & all defaults \\
\midrule
\inputrows{data/synth_win_rows.tex}
\bottomrule
\end{tabular}
\end{table}
\FloatBarrier

On ideal data, shortening the window removes the bias (\qty{0.15}{\second}:
\qty{0}{\percent}). With noise, a short window lets the noise of the differentiated $\dot u$
into the regressor, and the estimate falls \emph{below} the truth: this is the
errors-in-variables attenuation of \cref{ch:estimation}. The window is therefore a
bias--variance trade-off that has to be chosen against $\taueff$ and the noise level, and the
synthetic bench is where it can be tuned. The lag bias, finally, does not go away with any window.

\paragraph{Compensating the delay.} \Cref{tab:synth-comp} fits the lagged bag again with a copy
of the truth model whose \code{servo_delay_s} is set.

\begin{table}[htbp]
\centering
\caption{Total inertia $\mtot$ returned by the equation-error fit on the bag with a
\qty{0.15}{\second} actuation lag, for three values of \code{servo_delay_s} and two window
lengths (truth \qty{19.86}{\kilogram}).}
\label{tab:synth-comp}
\small
\begin{tabular}{@{}c l l@{}}
\toprule
\code{servo_delay_s} / \unit{\second} & window \qty{1}{\second} & window \qty{0.25}{\second} \\
\midrule
\inputrows{data/synth_comp_rows.tex}
\bottomrule
\end{tabular}
\end{table}
\FloatBarrier

With $-0.15$ the fit returns the same values as on the ideal bag (compare
\cref{tab:synth-window}), so the lag is fully removed and only the window bias remains. With
$+0.15$ the estimate never improves: it worsens to \qty{99}{\percent} with the
\qty{1}{\second} window and stays at \qty{25}{\percent} with \qty{0.25}{\second}. This is the
same sign convention as in \cref{tab:synth-prop}, now confirmed through a second, independent
code path.

\section{A bug found on the way}

The first attempt to run \file{propagate_model.py} on a synthetic bag crashed with
\code{Type 'std_msgs/msg/Float64' is unknown}. \code{write_augmented_bag()} in
\file{bag_io.py} built the new \code{/sysid/*} messages from the \emph{input} bag's
typestore, which only knows the types that happen to be recorded in it. The large autonomy
bags contain \code{Float64}, \code{NavSatFix} and the others by coincidence. A slim recording
made only for system identification would not. The fix registers every standard ROS~2 Jazzy
type the input bag does not define, and never overrides a recorded definition. A regression run on
\code{zermatt_rectangle_07} reproduces the committed summary digit for digit.

\section{Limitations}
\label{sec:synth-limits}

The bench imitates what the surge model assumes. It does not test effects outside it:
\begin{itemize}[nosep]
  \item surge only, at a constant heading. There is no yaw, sway, heave or pitch motion, and so
        no coupling;
  \item no current dynamics, waves or tether. Only a constant force offset is available;
  \item the EKF is a first-order lag, not a Kalman filter, and all sensor noise is white and Gaussian;
  \item the IMU sits at the centre of mass (no lever-arm accelerations) with a constant pitch;
  \item the motor is a pure delay plus a first-order lag, and the true thrust is the
        datasheet table times a constant;
  \item (added 16 October 2026) the thrust follows the \emph{logged} PWM, so the edges always
        coincide with a log sample. The effect of a slow \code{servo_output_raw} rate is therefore
        invisible on the bench. It is studied separately in \cref{tab:bias-servo};
  \item (added 16 October 2026) the IMU has perfect timestamps and no vibration, unlike the real
        \code{/imu/data} (\cref{sec:bias-study}).
\end{itemize}
The default imperfections are educated guesses. Once the delays and noise levels have been measured
at the lake (\cref{ch:timing}), they should replace the defaults so that the bench imitates
POLARIS rather than a generic vehicle.

\begin{incode}
\file{tests/test_synth_bag.py} (8 tests) checks that the ideal bag is reproduced by the
propagation to \qty{1e-6}{\meter\per\second}; that a lag $L$ is removed by
$\code{servo_delay_s} = -L$ and made worse by $+L$; that the events bracket the PWM edges; that
voltage and frames read back as from a real bag; that the IMU reads $-g\sin\theta_0$ at rest;
that the DVL velocity rotates back to the body velocity through \code{/tf_static}; and that the
\file{propagate_model.py} command line runs end to end on a synthetic bag.
\end{incode}

\begin{questions}
\begin{enumerate}
  \item Why must a fit be tested on the ideal bag before the bags with imperfections?
  \item The IMU shows the acceleration edge \qty{120}{\milli\second} after the logged PWM step. What
        value goes into \code{servo_delay_s}?
  \item In \cref{tab:synth-fit}, $\mtot$ of the thrust-scale bag is almost exactly right. Why is
        that not reassuring?
  \item Why does a shorter smoothing window \emph{under}estimate $\mtot$ on the noisy bag?
\end{enumerate}
\end{questions}
